% !TEX root = ../main.tex
\chapter{Axioms of set theory}\label{ch:l02}
\begin{paracol}{2}

\begin{sources}
\begin{tabularx}{\linewidth}{@{}l l L@{}}
\toprule
\textbf{Video} & \textbf{Length} & \textbf{Content}\\
\midrule
\seg{2}{1}{Lecture 2: Axioms of set theory} & 1:51:56 & The element relation; $\notin$, $\subseteq$, $=$;
the nine Zermelo--Fraenkel axioms with choice (ZFC): $\in$, empty set (with a fully formal proof of its
uniqueness), pair sets, union sets, replacement (restricted comprehension, complements), power sets,
infinity (the non-negative integers), choice, foundation. From 1:35:51 the video continues with the beginning of
the next session (choice and foundation).\\
\bottomrule
\end{tabularx}

\smallskip
\textbf{Transcript} \file{materials/transcripts/C02-V01.txt} (local whisper transcription: the video has no
YouTube captions).
\end{sources}

\begin{reading}[Where this lecture sits in the other notes]
Not treated in the other Schuller courses, which take sets for granted.
\end{reading}

% ==================================================================
\section{The element relation}\label{sec:se-in}

Last time propositional and predicate logic were introduced as the language in which set theory is to be
formulated; today that is what we do \ts{2}{1}{0:08}. Axiomatic set theory is built on one postulate: there is a
fundamental \emph{relation} called $\in$ \ts{2}{1}{0:08}.

\begin{definition}[New concepts: relation]\label{def:se-relation}
A \textbf{relation} is a predicate of two variables; ``relation'' is just shorter to say \ts{2}{1}{1:03}.
\end{definition}

Obviously $\in$ will be the element symbol we expect in set theory, but no preconceived notions: it is simply a
relation, called $\in$ \ts{2}{1}{1:28}. There will be \emph{no definition}, in the strict sense, of what $\in$ is,
apart from being a predicate of two variables, nor of what a set is \ts{2}{1}{1:53}\ts{2}{1}{2:16}. Instead we write
down nine axioms that speak of $\in$ and of sets; they teach us how to use $\in$ and what constitutes a set. Only
the interplay between $\in$ and what we call sets defines both \ts{2}{1}{2:38}\ts{2}{1}{3:03}. As said last time, this
is necessary if we want to write down mathematics from scratch, without any prior notions in terms of which a set
could be defined \ts{2}{1}{3:25}.

\paragraph{Overview of the axioms.} Schuller memorises them ``like a poem'' by their initial letters
\ts{2}{1}{3:46}\ts{2}{1}{4:11}:
\[
  \in\ \ \emptyset\qquad P\ \ U\ \ R\ \ P\qquad I\ \ C\ \ F
\]
\begin{itemize}
  \item $\in$, $\emptyset$: \textbf{basic existence axioms}: a property of the element relation, and the existence of
        an empty set \ts{2}{1}{4:33};
  \item P, U, R, P: heavily used \textbf{construction axioms}, which say how to build new sets from given ones
        (a recurring theme: given a structure, how to build new structures of the same type): \textbf{p}air sets,
        \textbf{u}nion sets, \textbf{r}eplacement (``a powerful thing''), \textbf{p}ower sets \ts{2}{1}{4:55}\ts{2}{1}{5:18};
  \item I, C: \textbf{further existence axioms} (infinity, choice), which could also be called more advanced
        construction axioms, and are historically newer \ts{2}{1}{5:41};
  \item F: the axiom of \textbf{foundation}, rather a \emph{non-existence} axiom: ``don't use sets of the following
        kind'' \ts{2}{1}{6:04}.
\end{itemize}
They are discussed in this order, partly because later axioms rely on earlier ones, partly to give an immediate
application of each \ts{2}{1}{6:30}\ts{2}{1}{6:53}.
\begin{coursenote}[The mnemonic]
The transcript has ``E E perp I C F''. Read as ``$\in$, $\emptyset$, PURP, ICF'': the two Es are the element relation
and the empty set, ``perp'' is P-U-R-P.
\end{coursenote}

\paragraph{Relations defined from $\in$.} A relation eats two variables; by tradition its symbol is written between
them, which saves two brackets and a comma \ts{2}{1}{7:39}\ts{2}{1}{8:02}.

\begin{definition}[New concepts: $\notin$, subset, equality of sets]\label{def:se-subset}
\begin{align*}
  x\notin y &:\Leftrightarrow \neg(x\in y), \\
  x\subseteq y &:\Leftrightarrow \forall a\colon (a\in x\Rightarrow a\in y), \\
  x=y &:\Leftrightarrow (x\subseteq y)\wedge(y\subseteq x)
\end{align*}
\ts{2}{1}{8:24}\ts{2}{1}{8:48}\ts{2}{1}{9:12}.
\end{definition}
So the equality sign, which one would normally take without questioning, is \emph{defined} in terms of the element
relation \ts{2}{1}{9:12}\ts{2}{1}{9:34}.

% ==================================================================
\section{The axiom on the element relation}\label{sec:se-ax-in}

The Zermelo--Fraenkel axioms of set theory \ts{2}{1}{9:57}. The first declares the relation between $\in$ and sets
\ts{2}{1}{10:21}.

\begin{definition}[Axiom on the $\in$-relation]\label{ax:se-in}
$x\in y$ is a proposition if and only if $x$ and $y$ are both sets \ts{2}{1}{10:48}\ts{2}{1}{11:09}. Formally:
$\forall x\,\forall y\colon (x\in y)\veebar\neg(x\in y)$ (\emph{tertium non datur}) \ts{2}{1}{23:04}\ts{2}{1}{23:33}.
\end{definition}
In predicate logic the nature of the variables was left entirely open; this axiom says that $x\in y$ is a
proposition, and $\in$ a relation, only on what we are going to call sets. It is the first thing we know about sets
\ts{2}{1}{11:37}. (The formal version is ``a little bit unnecessary'': we know what it means to be a proposition
\ts{2}{1}{23:59}.)

This seemed so trivial that for a long time it was not considered much of a condition; but it is. A famous
counterexample of something that is \emph{not} a set \ts{2}{1}{12:04}:

\begin{proposition}[Russell]\label{prop:se-russell}
Assume there is an object $u$ that contains all sets that do not contain themselves:
$\exists u\ \forall z\colon (z\in u\Leftrightarrow z\notin z)$, with $z$ ranging over sets
\ts{2}{1}{12:32}\ts{2}{1}{13:35}\ts{2}{1}{13:58}. Then $u$ is not a set.
\end{proposition}
\begin{proof}
If $u$ were a set, then by \cref{ax:se-in} $u\in u$ would be a proposition, true or false \ts{2}{1}{15:39}\ts{2}{1}{16:09}.
If $u\in u$ is true, $u$ is one particular $z$, so by the definition of $u$, $u\notin u$: contradiction; so the
assumption cannot be true \ts{2}{1}{16:33}\ts{2}{1}{17:03}. Then it must be false, since it is a proposition: $u\notin u$.
But $u$ contains all sets that do not contain themselves, so $u\in u$: contradiction again \ts{2}{1}{17:29}\ts{2}{1}{17:54}.
Both cases are contradictory, so $u$ is not a set \ts{2}{1}{18:17}.
\end{proof}
One may be disturbed by a set being, or not being, an element of itself, having been trained to distinguish sets
from their elements; the axiom of foundation will eventually say ``don't even look at such stuff''. But at this
point it is perfectly fine: if $z$ is a set, $z\in z$ is a proposition \ts{2}{1}{14:21}\ts{2}{1}{14:49}\ts{2}{1}{15:15}.
This is the \textbf{Russell paradox} (``Russell drama'') \ts{2}{1}{18:43}.

In \emph{naive set theory}, taught at school and in undergraduate courses, one easily writes
$u=\{z\mid z\notin z\}$, conditioning the elements of a set on a property, and so constructs objects that are not
sets \ts{2}{1}{19:07}\ts{2}{1}{19:38}\ts{2}{1}{20:02}. The contradictions of naive set theory were finally resolved
``in a big struggle by a number of people''; today we honour Zermelo and Fraenkel in particular. The first axiom
already shows the problem; eight more are needed not to run into difficulties \ts{2}{1}{20:31}\ts{2}{1}{21:00}.

\begin{history}[Russell's drama]
As told in class \ts{2}{1}{18:43}\ts{2}{1}{20:31}: Russell and Whitehead, Cambridge philosophers, logicians and
mathematicians, wrote the thick \emph{Principia Mathematica} on logic as the basic language of mathematics, their
life's work, built on naive set theory; when Russell discovered the contradiction he could not sleep, and had
nightmares of people in the Cambridge library throwing the volumes out of the window.\par
\smallskip
\emph{Checked.} Russell found the paradox in 1901 and wrote it to Frege in June 1902, as the second volume of
Frege's \emph{Grundgesetze der Arithmetik} (whose system it destroyed) was in press. \emph{Principia Mathematica}
(three volumes, 1910--1913) came \emph{after}, and avoids the paradox with the theory of types. The nightmare is
Russell's dream, told by G.~H.~Hardy in \emph{A Mathematician's Apology} (1940), of a librarian in the year 2100
deciding whether to throw away the last copy of \emph{Principia}. Zermelo's axioms: 1908; Fraenkel's replacement: 1922.
\end{history}

% ==================================================================
\section{The empty set}\label{sec:se-empty}

\begin{definition}[Axiom on the existence of an empty set]\label{ax:se-empty}
There exists a set that contains no elements \ts{2}{1}{21:57}. Formally: $\exists x\ \forall y\colon y\notin x$
\ts{2}{1}{22:29}.
\end{definition}

\begin{theorem}\label{thm:se-empty-unique}
There is only one empty set; it is therefore given a name, $\emptyset$ \ts{2}{1}{24:25}.
\end{theorem}
\begin{proof}[Standard textbook-style proof]
Let $x$ and $x'$ be empty sets; we show $x=x'$ \ts{2}{1}{24:59}\ts{2}{1}{25:38}. For any $y$, $y\in x$ is false, since
$x$ has no elements. So the implication $y\in x\Rightarrow y\in x'$ is true, whatever is written after the arrow (ex
falso quodlibet), independently of $y$ \ts{2}{1}{26:11}. By the definition of $\forall$,
$\forall y\colon(y\in x\Rightarrow y\in x')$ is true, i.e.\ $x\subseteq x'$ \ts{2}{1}{26:39}. Conversely, $y\in x'$ is
always false, so $\forall y\colon (y\in x'\Rightarrow y\in x)$, i.e.\ $x'\subseteq x$ \ts{2}{1}{27:20}\ts{2}{1}{27:47}.
Together, $x=x'$ \ts{2}{1}{28:13}.
\end{proof}

Didn't we define what a proof is? This textbook proof does not have that form \ts{2}{1}{28:37}. So here, once and only
once in the course, is the formal version \ts{2}{1}{28:37}\ts{2}{1}{29:49}. The assumptions become the axioms of the system
(propositional logic alone would have an empty list) \ts{2}{1}{30:57}\ts{2}{1}{31:19}:
\[
  a_1:\ \forall y\colon y\notin x,\qquad a_2:\ \forall y\colon y\notin x'.
\]
We need steps $q_1,\dots,q_M$ with $q_M$ ``$x=x'$'', each justified by (T), (A) or (M) of \cref{def:lo-proof}
\ts{2}{1}{30:13}\ts{2}{1}{31:41}:
\begin{enumerate}
  \item[$q_1$] $\bigl(\forall y\colon y\notin x\bigr)\Rightarrow\bigl(\forall y\colon (y\in x\Rightarrow y\in x')\bigr)$
        \hfill (T): a tautology, by ex falso quodlibet \ts{2}{1}{32:04}\ts{2}{1}{32:49};
  \item[$q_2$] $\forall y\colon y\notin x$ \hfill (A): axiom $a_1$ \ts{2}{1}{33:12}\ts{2}{1}{33:34};
  \item[$q_3$] $\forall y\colon (y\in x\Rightarrow y\in x')$, i.e.\ $x\subseteq x'$ \hfill (M) with $q_1,q_2$:
        $(q_1\wedge q_2)\Rightarrow q_3$ is a tautology \ts{2}{1}{33:56}\ts{2}{1}{34:40}\ts{2}{1}{35:24};
  \item[$q_4$] $\bigl(\forall y\colon y\notin x'\bigr)\Rightarrow\bigl(\forall y\colon (y\in x'\Rightarrow y\in x)\bigr)$
        \hfill (T) \ts{2}{1}{36:08};
  \item[$q_5$] $\forall y\colon y\notin x'$ \hfill (A): axiom $a_2$ \ts{2}{1}{36:30};
  \item[$q_6$] $\forall y\colon (y\in x'\Rightarrow y\in x)$, i.e.\ $x'\subseteq x$ \hfill (M) with $q_4,q_5$
        \ts{2}{1}{36:52}\ts{2}{1}{37:37};
  \item[$q_7$] $x=x'$ \hfill (M) with $q_3,q_6$ \ts{2}{1}{37:59}.
\end{enumerate}
\begin{coursenote}[Which steps the modus ponens uses]
On the board Schuller first said that $q_6$ used steps 5 and 6, then corrected himself: 4 and 5
\ts{2}{1}{37:37}. Strictly, propositions are not \emph{equal} but \emph{equivalent}, so the steps should be written with
$\Leftrightarrow$ rather than $=$ if one is picky \ts{2}{1}{30:35}.
\end{coursenote}
``Something you probably want to have a look at again at home'' \ts{2}{1}{35:01}. The formal proof is nothing more
than the textbook proof broken down pedantically, but now we are certain it is a proof \ts{2}{1}{38:21}. This is the
spirit of axiomatic set theory: two levels, the usual language and the precise first-order version that mirrors it
\ts{2}{1}{38:43}\ts{2}{1}{39:05}. We will not follow the strict technique further: proofs get longer and longer, and no
book spells out all its proofs this way; but it must be possible in principle \ts{2}{1}{39:27}\ts{2}{1}{39:50}.

A student asks whether $x\subseteq y$ and $x\in y$ imply one another. No, in neither direction; finding
counterexamples for both is easy \ts{2}{1}{29:02}\ts{2}{1}{29:25}.
\begin{exercise}[$\subseteq$ versus $\in$]\label{ex:se-subset-in}
Find sets with $x\subseteq y$ but $x\notin y$, and with $x\in y$ but $x\not\subseteq y$ (use $\emptyset$ and the sets of
the next sections).
\end{exercise}

% ==================================================================
\section{Pair sets}\label{sec:se-pair}

``A bit tiring'' to go through all nine axioms, but this is the foundation of all of mathematics. A field such as
$\R$ also needs nine axioms (commutativity, associativity, neutral and inverse element for $+$, the same for
$\cdot$, one distributive law), so building all of mathematics on nine axioms is not too bad
\ts{2}{1}{40:14}\ts{2}{1}{40:36}\ts{2}{1}{40:58}.

\begin{definition}[Axiom on pair sets]\label{ax:se-pair}
Let $x$ and $y$ be sets. Then there exists a set that contains as its elements precisely $x$ and $y$
\ts{2}{1}{40:58}\ts{2}{1}{41:20}. Formally: $\forall x\,\forall y\ \exists m\ \forall u\colon\bigl(u\in m\Leftrightarrow
(u=x\vee u=y)\bigr)$ \ts{2}{1}{42:09}\ts{2}{1}{42:32}. Notation: $m=\{x,y\}$ \ts{2}{1}{42:56}\ts{2}{1}{43:23}.
\end{definition}
The notation is just a symbol (``I could put a little flower in the upper right corner as well''); what the axiom
ensures is that, if $x$ and $y$ are sets, this object \emph{is a set} \ts{2}{1}{43:23}\ts{2}{1}{43:46}.

\paragraph{Does the order matter?} Is $\{x,y\}=\{y,x\}$? Yes: $a\in\{x,y\}\Rightarrow a\in\{y,x\}$ and conversely, so
each is a subset of the other \ts{2}{1}{44:09}\ts{2}{1}{44:30}\ts{2}{1}{44:52}. ``We don't agree on it, we prove it'';
that the whole thing is a set is guaranteed by the axiom \ts{2}{1}{45:13}.

\begin{definition}[New concepts: one-element set]\label{def:se-singleton}
For a set $x$, $\{x\}:=\{x,x\}$, a set by the pair set axiom \ts{2}{1}{45:13}\ts{2}{1}{45:38}. Since equality is defined
through $\in$, writing $x$ twice adds no information \ts{2}{1}{46:00}.
\end{definition}
So far we have one- and two-element sets, nothing more: at this stage there is no justification for writing
$\{x,y,z\}$ \ts{2}{1}{46:22}. The most demanding task now is to get rid of old prejudices about sets, and to see how
what we knew follows from the axioms (some things will not be allowed), much as one learns that commutativity of
the real numbers is something special that must be stated as a law \ts{2}{1}{46:44}\ts{2}{1}{47:08}\ts{2}{1}{47:33}.

% ==================================================================
\section{Union sets}\label{sec:se-union}

\begin{definition}[Axiom on union sets]\label{ax:se-union}
Let $x$ be a set. Then there exists a set whose elements are precisely the elements of the elements of $x$
\ts{2}{1}{47:57}\ts{2}{1}{48:20}\ts{2}{1}{48:46}. Notation: $\bigcup x$ \ts{2}{1}{49:58}\ts{2}{1}{50:22}.
\end{definition}

\begin{exercise}[Axioms in first-order language]\label{ex:se-formal}
Homework for every further axiom: convert the verbal statement into first-order predicate logic, as done above for
the empty set and the pair sets: ``if you think you understand it, put it in first-order predicate logic''
\ts{2}{1}{49:13}\ts{2}{1}{49:35}.
\end{exercise}

Don't we need at least two sets to build a union? No: the set $x$ already collects all the sets to be put into the
union \ts{2}{1}{49:35}\ts{2}{1}{49:58}.

\begin{example}[Union of two one-element sets]
Let $a,b$ be sets. $\{a\}$ and $\{b\}$ are sets (pair set axiom), hence so is $x=\{\{a\},\{b\}\}$ (pair set axiom
again) \ts{2}{1}{52:14}. Its union contains the elements of the elements of $x$: $\bigcup x=\{a,b\}$. Here we already
knew this is a set; the union set axiom says the result is a set whatever $x$ \ts{2}{1}{52:39}\ts{2}{1}{53:03}.
\end{example}

\begin{example}[A three-element set]
For sets $a,b,c$, $x=\{\{a,b\},\{c\}\}$ is a set, so $\bigcup x$ is a set by the union set axiom; it deserves a name:
$\{a,b,c\}:=\bigcup x$ \ts{2}{1}{53:03}\ts{2}{1}{53:29}\ts{2}{1}{53:56}. ``In elementary school you learn this is a set,
but how do you know? It's just claimed. Now the axiom tells you'' \ts{2}{1}{54:21}.
\end{example}

\begin{definition}[New concepts: finite sets with specified elements]\label{def:se-finite}
Let $a_1,\dots,a_N$ be sets (the indices are again pre-mathematical stroke numbers) \ts{2}{1}{54:45}. For $N\ge3$ define
recursively
\[
  \{a_1,\dots,a_N\}:=\bigcup\bigl\{\{a_1,\dots,a_{N-1}\},\{a_N\}\bigr\},
\]
the recursion being started at $N=1,2$ by the pair set axiom \ts{2}{1}{55:27}\ts{2}{1}{55:56}.
\end{definition}
This gives everything we would informally call a finite set with specified elements \ts{2}{1}{55:56}.

\begin{intuition}[As many sets as fit into a set]
The big thing about the axiom is the assumption that the sets to be united are collected \emph{in a set} $x$.
Russell's $u$ is not a set, so it can never stand in place of $x$: one can only build unions of ``as many sets as
still fit into a set'' \ts{2}{1}{56:20}\ts{2}{1}{56:52}\ts{2}{1}{57:19}. This is much more than taking finitely many
sets (it contains, roughly, taking infinitely many), but not more sets than fit into a set: subtle, but very
precise \ts{2}{1}{57:44}. The elementary-school notion ``a set is a collection of elements'' comes back, now with a
justification from the first four axioms \ts{2}{1}{57:44}\ts{2}{1}{58:09}.
\end{intuition}

% ==================================================================
\section{Replacement and restricted comprehension}\label{sec:se-replacement}

Pedagogically, replacement comes before power sets, although it is not logically necessary
\ts{2}{1}{58:32}\ts{2}{1}{59:03}. ``A very innocent-sounding name; one of the most powerful axioms'' \ts{2}{1}{59:41}.

\begin{definition}[New concepts: functional relation, image]\label{def:se-functional}
A relation $R$ is \textbf{functional} if for every $x$ there exists precisely one $y$ such that $R(x,y)$
\ts{2}{1}{1:01:03}\ts{2}{1}{1:01:38}. The \textbf{image} of a set $m$ under a functional relation $R$ ``consists of'' all
those $y$ for which there is an $x\in m$ with $R(x,y)$ \ts{2}{1}{1:02:52}\ts{2}{1}{1:03:15}\ts{2}{1}{1:03:37}.
\end{definition}
This is what one calls a function: to every $x$ precisely one $y$ is assigned (to another $x$ maybe another $y$) but
at this stage there is no talk of $x$ or $y$ being sets: it is just a relation \ts{2}{1}{1:02:03}\ts{2}{1}{1:02:29}.
And ``your alarm bell should ring'' at the harmless word \emph{consists}: who says that all those $y$ are the
elements of a set \ts{2}{1}{1:03:58}\ts{2}{1}{1:04:19}?

\begin{exercise}[``There exists precisely one'']\label{ex:se-exists-one}
Define $\exists!\,y\colon P(y)$ (``there exists precisely one $y$'') in terms of $\forall$, $\exists$ and $=$
(homework in class, ``instead of an examples class'') \ts{2}{1}{1:01:38}\ts{2}{1}{1:02:03}.
\end{exercise}

\begin{definition}[Axiom of replacement]\label{ax:se-replacement}
Let $R$ be a functional relation and $m$ a set. Then the image of $m$ under $R$ is again a set
\ts{2}{1}{1:00:10}\ts{2}{1}{1:00:39}.
\end{definition}
A construction axiom: start with a set, construct the image under a functional relation, and the axiom guarantees
the result is a set \ts{2}{1}{1:00:39}\ts{2}{1}{1:04:42}. You have met it already, even at school, but in a weaker form,
restricted comprehension, which follows from it. The full axiom is needed to cover all of modern mathematics, but
in this course only its weaker consequence will be used \ts{2}{1}{1:05:03}\ts{2}{1}{1:05:23}\ts{2}{1}{1:05:45}.

\begin{theorem}[Principle of restricted comprehension]\label{thm:se-comprehension}
Let $P$ be a predicate of one variable and $m$ a set. Then the elements $y\in m$ for which $P(y)$ holds constitute a
set, denoted $\{y\in m\mid P(y)\}$ \ts{2}{1}{1:06:20}\ts{2}{1}{1:07:28}\ts{2}{1}{1:07:52}.
\end{theorem}
\begin{proof}
Two cases \ts{2}{1}{1:11:10}.
\begin{enumerate}
  \item There is no $y\in m$ with $P(y)$. Then the symbol is defined as $\emptyset$, which is a set; replacement is not
        needed \ts{2}{1}{1:11:39}\ts{2}{1}{1:12:42}.
  \item There is a $y\in m$ with $P(y)$; pick one, $\hat y$ \ts{2}{1}{1:13:19}. Define the relation
        \[
          R(x,y):\Leftrightarrow\bigl(P(x)\wedge y=x\bigr)\vee\bigl(\neg P(x)\wedge y=\hat y\bigr)
        \]
        \ts{2}{1}{1:13:49}\ts{2}{1}{1:14:32}\ts{2}{1}{1:14:52}. It is functional: if $P(x)$ holds, $x$ is sent to itself;
        if not, to the one fixed $\hat y$ \ts{2}{1}{1:15:15}\ts{2}{1}{1:15:36}. Define $\{y\in m\mid P(y)\}$ as the image of
        $m$ under $R$, a set by replacement \ts{2}{1}{1:15:59}\ts{2}{1}{1:16:21}.
\end{enumerate}
``You need to stare at this at home for five minutes'' \ts{2}{1}{1:15:59}\ts{2}{1}{1:16:43}.
\end{proof}
Replacement implies restricted comprehension, but not conversely \ts{2}{1}{1:06:20}.

\begin{remark}[Restricted quantifiers]\label{rem:se-restricted-q}
The proof uses quantifiers restricted to a set, which first-order logic does not have; they are short forms
\ts{2}{1}{1:11:39}\ts{2}{1}{1:12:12}:
\[
  \forall y\in m\colon P(y):\Leftrightarrow\forall y\colon(y\in m\Rightarrow P(y)),\qquad
  \exists y\in m\colon P(y):\Leftrightarrow\neg\bigl(\forall y\in m\colon\neg P(y)\bigr),
\]
so that $\exists y\in m\colon P(y)\Leftrightarrow\exists y\colon(y\in m\wedge P(y))$. Schuller first wrote the
existential version with $\Rightarrow$; a student objected, and the correct form is the one above: for $\forall$ an
implication, so that a $y$ outside $m$ does not affect the truth of the statement, and $\exists$ derived from it by
the double negation, which turns the implication into a conjunction \ts{2}{1}{1:20:07}\ts{2}{1}{1:20:47}\ts{2}{1}{1:21:26}\ts{2}{1}{1:21:57}.
``Very good observation. The laces of my shoes came undone because of this mistake'' \ts{2}{1}{1:22:28}.
\end{remark}

\begin{intuition}[Restricted versus universal comprehension]
Not to be confused with the inconsistent \textbf{principle of universal comprehension}, Russell's problem: ``the $y$
for which $P(y)$ holds form a set'' \ts{2}{1}{1:08:48}\ts{2}{1}{1:09:27}. That could be all $y$ with $y\notin y$
\ts{2}{1}{1:10:05}. The difference is that a set $m$ must be given from the start: the selection can only kick out some
elements of $m$, never make anything bigger, so it is not unbounded. The ``restricted'' is exactly ``$\in m$''
\ts{2}{1}{1:10:05}\ts{2}{1}{1:10:39}\ts{2}{1}{1:11:10}.
\end{intuition}

\begin{definition}[New concepts: complement]\label{def:se-complement}
For sets $u\subseteq m$, $m\setminus u:=\{x\in m\mid x\notin u\}$, a set by restricted comprehension, ultimately by
replacement \ts{2}{1}{1:17:27}\ts{2}{1}{1:17:49}\ts{2}{1}{1:18:11}.
\end{definition}

\begin{exercise}[Intersections]\label{ex:se-intersection}
Let $x$ be a set. Define the intersection of all the elements of $x$, using restricted comprehension (homework in
class) \ts{2}{1}{1:18:37}\ts{2}{1}{1:18:58}. Note the asymmetry: the existence of unions had to be postulated; intersections
follow from replacement \ts{2}{1}{1:19:22}.
\end{exercise}

% ==================================================================
\section{Power sets}\label{sec:se-power}

The power set is the collection of all subsets of a set \ts{2}{1}{1:22:28}\ts{2}{1}{1:23:36}. Historically, in naive set
theory, the principle of universal comprehension was thought to be needed to define it: ``all the $y$ that are
subsets of $m$'' \ts{2}{1}{1:23:36}\ts{2}{1}{1:24:28}. Given only $m$, we cannot say from which bigger set these subsets
come, so restricted comprehension does not help \ts{2}{1}{1:24:56}. And one wants to apply power sets repeatedly,
the power set of the power set\dots, going through the whole universe of sets; to \emph{generate} sets one cannot know
beforehand where they come from \ts{2}{1}{1:24:56}\ts{2}{1}{1:25:22}. There is no other choice than to postulate
\ts{2}{1}{1:25:53}:

\begin{definition}[Axiom on the existence of power sets]\label{ax:se-power}
Let $m$ be a set. Then there exists a set, denoted $\mathcal P(m)$, whose elements are precisely the subsets of $m$
\ts{2}{1}{1:25:53}\ts{2}{1}{1:26:37}.
\end{definition}

\begin{example}[Power set of a two-element set]
For $m=\{a,b\}$: $\mathcal P(m)=\bigl\{\emptyset,\{a\},\{b\},\{a,b\}\bigr\}$: the empty set (a subset of any set), the
one-element subsets, and $m$ itself \ts{2}{1}{1:26:37}\ts{2}{1}{1:27:37}.
\end{example}
Easy to explain how to build it; but is the result a set? Neither the non-existent universal principle nor
restricted comprehension can say so: an extra axiom is needed \ts{2}{1}{1:28:04}\ts{2}{1}{1:28:29}. ``Not very
economical'', but the power sets we want \ts{2}{1}{1:28:29}.

% ==================================================================
\section{Infinity}\label{sec:se-infinity}

\begin{definition}[Axiom of infinity]\label{ax:se-infinity}
There exists a set that contains the empty set and, with every one of its elements $y$, also contains $\{y\}$
\ts{2}{1}{1:29:23}\ts{2}{1}{1:29:48}.
\end{definition}
Without it we do not get standard mathematics \ts{2}{1}{1:29:23}. (``With every one of its elements'' refers to the
elements of the set being postulated, not of $\emptyset$, which has none \ts{2}{1}{1:29:48}.)

\begin{remark}[The simplest such set]\label{rem:se-naturals}
Informally, the simplest such set is $\bigl\{\emptyset,\{\emptyset\},\{\{\emptyset\}\},\{\{\{\emptyset\}\}\},\dots\bigr\}$
\ts{2}{1}{1:30:39}\ts{2}{1}{1:31:22}. Too many brackets for elementary school, so trivial names are introduced:
$0:=\emptyset$, $1:=\{\emptyset\}$, $2:=\{\{\emptyset\}\}$, and so on \ts{2}{1}{1:31:55}\ts{2}{1}{1:32:21}. The axiom guarantees
other sets too; this is one particular of them \ts{2}{1}{1:32:21}.
\end{remark}

\begin{corollary}\label{cor:se-N}
The non-negative integers form a set according to axiomatic set theory \ts{2}{1}{1:32:50}.
\end{corollary}
Without the axiom of infinity one could not prove it \ts{2}{1}{1:33:17}.
\begin{coursenote}[The precise construction]
``Informally speaking'': the set of \cref{rem:se-naturals} is obtained from any set given by the axiom, by
restricted comprehension, as the smallest one (the intersection of all subsets with the property). These are
Zermelo's numerals, $n+1=\{n\}$; the now more common von Neumann numerals use $n+1=n\cup\{n\}$.
\end{coursenote}

What next: do we require the real numbers to exist? As a \emph{set} (nothing further) they can be understood as, or
defined as, the power set of the integers \ts{2}{1}{1:33:17}\ts{2}{1}{1:33:47}. Then the power set of the reals is a bigger
set, its power set bigger still, and so on through the universe of sets \ts{2}{1}{1:33:47}\ts{2}{1}{1:34:19}. And the only set
whose existence was explicitly demanded so far is the empty set: from it we already arrived at the reals
\ts{2}{1}{1:34:19}.
\begin{coursenote}[``The reals are the power set of the integers'']
As a statement about sets this means: $\mathcal P(\mathbb N)$ has the cardinality of $\R$, so it can serve as the
underlying set; $\R$ as the field of real numbers is built differently (e.g.\ Dedekind cuts of $\mathbb Q$).
\end{coursenote}

The two remaining axioms are announced: C, which one may put in or leave out (ZF with or without choice), and F,
which forbids ``bootstrapping'' sets that contain themselves; with all axioms, essentially every set of standard
mathematics is constructed from the empty set alone. ``So set theory is very simple: it's just the empty set''
\ts{2}{1}{1:34:48}\ts{2}{1}{1:35:19}.

% ==================================================================
\section{Choice}\label{sec:se-choice}

The video continues with the next session: ``welcome back''; seven axioms discussed, two to go
\ts{2}{1}{1:35:51}.

\begin{definition}[Axiom of choice]\label{ax:se-choice}
Let $x$ be a set whose elements are non-empty and mutually disjoint (the intersection of any two elements of $x$ is
$\emptyset$). Then there exists a set $y$ that contains exactly one element of each element of $x$
\ts{2}{1}{1:36:28}\ts{2}{1}{1:37:08}.
\end{definition}
Intuitively: from each of the elements of $x$, which are all sets, pick one element, and collect them all in $y$
\ts{2}{1}{1:37:08}\ts{2}{1}{1:38:12}. Sometimes $y$ is called a ``dark'' set: given $x$, there is no prescription of how to
pick the elements \ts{2}{1}{1:38:41}.

\begin{intuition}[Shoes and socks]
An intuition, ``I think by Russell'' \ts{2}{1}{1:38:41}\ts{2}{1}{1:39:14}: if $x$ consists of pairs of shoes,
$\{\{\text{left}_1,\text{right}_1\},\{\text{left}_2,\text{right}_2\},\dots\}$, no axiom of choice is needed: there is an
algorithm, ``always pick the left shoe'' \ts{2}{1}{1:39:57}. For pairs of socks, two different individual socks that
cannot be distinguished, one must invoke the axiom of choice: there is no further means of telling which one is
picked \ts{2}{1}{1:40:37}\ts{2}{1}{1:41:15}. It is just an analogy \ts{2}{1}{1:41:52}.
\end{intuition}

A student asks whether it should say ``one element of each \emph{set} of $x$''. The point is that everything that
appears is a set: ``element'' only emphasises that something lies in something else \ts{2}{1}{1:41:52}\ts{2}{1}{1:42:22}.

\begin{remark}[Independence and uses]\label{rem:se-choice}
The axiom of choice is independent of the other axioms: set theory can be done with or without it; standard
mathematics uses it, and so will we \ts{2}{1}{1:42:22}\ts{2}{1}{1:43:27}. Two theorems of this course need it
\ts{2}{1}{1:43:27}\ts{2}{1}{1:44:12}:
\begin{enumerate}
  \item every vector space has a basis;
  \item for every equivalence relation there exists a complete system of representatives (equivalence relations
        come later).
\end{enumerate}
\end{remark}

\begin{history}[Zermelo and the axiom of choice]
As said in class (``I think, I'm not entirely sure'') \ts{2}{1}{1:44:52}\ts{2}{1}{1:45:19}: the axiom was introduced by
Zermelo to prove a theorem that had long been unproven, so it was first suspected of being made up just for that;
today it is used throughout, and nothing about it is suspicious.\par
\smallskip
\emph{Checked.} E.~Zermelo, ``Beweis, da\ss{} jede Menge wohlgeordnet werden kann'', \emph{Math.\ Ann.}\ 59 (1904)
514--516: the well-ordering theorem, conjectured by Cantor. Independence: G\"odel (1938) showed choice consistent with
ZF, P.~Cohen (1963) that its negation is too. The shoes-and-socks image is indeed Russell's (\emph{Introduction to
Mathematical Philosophy}, 1919).
\end{history}

% ==================================================================
\section{Foundation}\label{sec:se-foundation}

The last axiom is not an existence or construction axiom but a non-existence axiom; notwithstanding its name it was
one of the last to be introduced \ts{2}{1}{1:45:19}\ts{2}{1}{1:46:05}.

\begin{definition}[Axiom of foundation]\label{ax:se-foundation}
Every non-empty set $x$ contains an element $y$ that has none of its elements in common with $x$:
$x\ne\emptyset\Rightarrow\exists y\in x\colon y\cap x=\emptyset$ \ts{2}{1}{1:46:31}\ts{2}{1}{1:47:01}.
\end{definition}

\begin{corollary}\label{cor:se-no-self}
There is no set $x$ with $x\in x$ \ts{2}{1}{1:47:25}.
\end{corollary}
\begin{proof}
(Not written out in class.) Let $x$ be a set; $\{x\}$ is a non-empty set, so by foundation it has an element with no
element in common with $\{x\}$. Its only element is $x$, so $x\cap\{x\}=\emptyset$, i.e.\ $x\notin x$.
\end{proof}
Recall the attempt to build the set of all sets that do not contain themselves: that is a non-set, since it violates
the first axiom; and self-containing sets could not be excluded by a principle of universal comprehension, since
there is none, only a restricted one. So the axiom of foundation has to take care of it
\ts{2}{1}{1:48:06}\ts{2}{1}{1:48:30}\ts{2}{1}{1:48:55}. It also excludes defining sets self-referentially; without such sets
every set must ultimately be built from the empty set, as the non-negative integers were by the axiom of infinity
\ts{2}{1}{1:49:20}\ts{2}{1}{1:49:50}.

\paragraph{Summary.} The nine axioms $\in,\emptyset$, P, U, R, P, I, C, F should be memorised; every time a set is
constructed, ask whether it can actually be constructed from them \ts{2}{1}{1:49:50}\ts{2}{1}{1:50:17}. Their totality is
\textbf{Zermelo--Fraenkel set theory with the axiom of choice}, \textbf{ZFC}; without C, ZF \ts{2}{1}{1:50:17}\ts{2}{1}{1:50:56}.
There are other, weaker axiomatisations (not really stronger ones, as far as Schuller knows); as far as anybody knows,
all of standard mathematics is built on ZFC and can, in principle, be constructed from it
\ts{2}{1}{1:50:56}\ts{2}{1}{1:51:19}.

% ==================================================================
\section*{Key concepts}
\begin{keyconcepts}
\textbf{Retrieval questions.}
\begin{enumerate}
  \item How are $\subseteq$ and $=$ defined from $\in$?
  \item Why is $\{z\mid z\notin z\}$ not a set, and which axiom shows it?
  \item List the nine axioms of ZFC with their character (existence, construction, non-existence).
  \item Why is $\{a,b,c\}$ a set? Why is $\{x\}$ one?
  \item State replacement and derive restricted comprehension from it. Why is universal comprehension inconsistent?
  \item Why must the existence of power sets be postulated? Which set does infinity give?
  \item Where will the axiom of choice be needed? What does foundation exclude?
\end{enumerate}
\textbf{Concepts.}
\begin{itemize}
  \item $x\subseteq y:\Leftrightarrow\forall a\colon(a\in x\Rightarrow a\in y)$; $x=y:\Leftrightarrow x\subseteq y\wedge y\subseteq x$ (\cref{def:se-subset}).
  \item Russell's paradox (\cref{prop:se-russell}); uniqueness of $\emptyset$ with its formal proof (\cref{thm:se-empty-unique}).
  \item Pair sets $\{x,y\}$, $\{x\}$; unions $\bigcup x$; finite sets $\{a_1,\dots,a_N\}$ (\cref{ax:se-pair,ax:se-union,def:se-finite}).
  \item Functional relation, image; replacement; $\{y\in m\mid P(y)\}$ (\cref{ax:se-replacement,thm:se-comprehension}).
  \item $\mathcal P(m)$; infinity and $0=\emptyset$, $1=\{\emptyset\}$, \dots (\cref{ax:se-power,ax:se-infinity}).
  \item Choice (\cref{ax:se-choice}); foundation, no $x\in x$ (\cref{ax:se-foundation,cor:se-no-self}); ZFC.
\end{itemize}
\end{keyconcepts}

\flushside
\end{paracol}
